%; whizzy chapter -dvi
% -initex iniptex -latex platex -format platex -bibtex jbibtex -fmt fmt
% 以上 whizzytex を使用する場合の設定。
 
%     Tokyo Debian Meeting resources
%     Copyright (C) 2012 Junichi Uekawa
%     Copyright (C) 2011 Nobuhiro Iwamatsu

%     This program is free software; you can redistribute it and/or modify
%     it under the terms of the GNU General Public License as published by
%     the Free Software Foundation; either version 2 of the License, or
%     (at your option) any later version.

%     This program is distributed in the hope that it will be useful,
%     but WITHOUT ANY WARRANTY; without even the implied warranty of
%     MERCHANTABILITY or FITNESS FOR A PARTICULAR PURPOSE.  See the
%     GNU General Public License for more details.

%     You should have received a copy of the GNU General Public License
%     along with this program; if not, write to the Free Software
%     Foundation, Inc., 51 Franklin St, Fifth Floor, Boston, MA  02110-1301 USA

%  preview (shell-command (concat "evince " (replace-regexp-in-string "tex$" "pdf"(buffer-file-name)) "&"))

%%ここからヘッダ開始。

\documentclass[mingoth,a4paper]{jsarticle}
\usepackage{monthlyreport}
% 日付を定義する、毎月変わります。
\newcommand{\debmtgyear}{2014}
\newcommand{\debmtgmonth}{10}
\newcommand{\debmtgdate}{25}
% started from zero:
% (let ((year 2013) (month 7)) (+ (* (- year 2005) 12) month -1))
\newcommand{\debmtgnumber}{119}

\begin{document}

\begin{titlepage}
\thispagestyle{empty}
% タイトルページ:編集必要な部分は最初のマクロに飛ばすこと

\vspace*{-2cm}
第\debmtgnumber{}回 東京エリア Debian 勉強会資料\\
\hspace*{-2cm}
\includegraphics{image2012-natsu/dotdeb.pdf}\\
\hfill{}\debmtgyear{}年\debmtgmonth{}月\debmtgdate{}日

% ここはアップデートすること
% 全角文字にしないとフォントのサイズが合わないので注意
\rotatebox{10}{\fontsize{30}{30} {\gt 特集：Debian x LibreOffice}}\\

\vspace*{-2cm}
\hfill{}\includegraphics[height=6cm]{image200502/openlogo-nd.eps}
\end{titlepage}

\newpage

\begin{minipage}[b]{0.2\hsize}
 \definecolor{titleback}{gray}{0.9}
 \colorbox{titleback}{\rotatebox{90}{\fontsize{80}{80} {\gt デビアン勉強会} }}
\end{minipage}
\begin{minipage}[b]{0.8\hsize}
\hrule
\vspace{2mm}
\hrule
\begin{multicols}{2}
\tableofcontents
\end{multicols}
\vspace{2mm}
\hrule
\end{minipage}

\dancersection{事前課題}{野島 貴英}

今回の事前課題は以下です:
\begin{enumerate}
 \item 本日、何の作業をやるかを宣言ください。
\end{enumerate}
この課題に対して提出いただいた内容は以下です。
\begin{multicols}{2}
{\small
\begin{prework}{ 野島　貴英 }
もちろん！Jessieインストーラテスト会に参加します。
よろしくお願いします。
\end{prework}

\begin{prework}{ henrich }
フリーズ前の課題洗い出しとか…。
\end{prework}

\begin{prework}{ dictoss(杉本　典充) }
Jessieインストーラテスト会に参加します。
\end{prework}

\begin{prework}{ 吉野(yy\_{}y\_{}ja\_{}jp) }
DDTSS
\end{prework}

\begin{prework}{ kenhys }
groonga-normalizer-mysqlのITP関連
\end{prework}

\begin{prework}{ Hasegawa }
Jessieインストーラテスト会に参加します。
よろしくお願いします。
\end{prework}

}
\end{multicols}

\dancersection{Debian Trivia Quiz}{野島 貴英}

 Debianの昨今の話題についてのQuizです。

今回の出題範囲は\url{debian-devel-announce@lists.debian.org} や \url{debian-news@lists.debian.org}に投稿された
内容などからです。

\begin{multicols}{2}
%; whizzy-master ../debianmeetingresume201311.tex
% 以上の設定をしているため、このファイルで M-x whizzytex すると、whizzytexが利用できます。
%

\santaku
{2014/9/14時点でJessieでサポートされると残念ながら「言われなかった」アーキテクチャはどれ}
{amd64}
{powerpc}
{sparc}
{C}
{2014/9/14時点でサポートされると言われたアーキテクチャは：amd64,armel,armhf,i386,kfreebsd-amd64/kfreebsd-i386/mips/mipsel/powerpc/s390x。なお、arm64/ppc64elは好調な頑張りとのことで、この調子が続けば入るかも？？という状況。}

\santaku
{2014/11/1時点でJessieでサポートされるかどうかについて依然として懸念のあるものはどれ}
{hurd}
{kFreeBSD}
{s390x}
{B}
{kFreeBSDは頑張って欲しいとのこと。最終ジャッジは2014/11/1に行われる。なお、s390xは9/14では懸念なし、hurdはすでにリリースは無理との判断になっている。}

\santaku
{2014/9/4にtesting入りしたデスクトップ環境は以下のどれ}
{xfce}
{cinnamon}
{unity}
{B}
{cinnamonはGTK+3を利用して作られたデスクトップ環境です。元々はLinux MINT向けにGNOME Shellからforkしたデスクトップシェルでした。開発が続き、デスクトップシェルから発展して遂にデスクトップ環境となりました。}

\santaku
{2014/10/18にリリースされたwheezyのDebianのバージョンはいくら？}
{7.7}
{7.8}
{7.9}
{A}
{7回目のアップデートとなります。早速アップグレードしましょう！新規にインストールする安定版ならDebian7.7からやりましょう!}

\santaku
{OPWの呼びかけがdebian-devel-announceで2014/9/30に行われました。ところでOPWって何の略？}
{One-time PassWord}
{One-Piece-Woven technology}
{Outreach-Program-for-Women}
{C}
{Outreach-Program-for-Women(略してOPW)は、GNOME Foundationが始めた、FOSSのプロジェクトの貢献者にもっと女性を増やそう！という活動となります。FOSSの貢献者の男女構成比は、圧倒的に男の割合が高いといういびつな傾向があるため、こちらを是正しようとする活動です。}

\santaku
{2014/10/14にサービスを近いうちに閉じるよとアナウンスのあったサイトはどれでしょう？}
{githubredir.debian.net}
{tracker.debian.org}
{rtc.debian.org}
{A}
{githubredir.debian.netは、github上にあるupstreamのソースについて、リリースの為にバージョン番号でtag打ちされたバージョンのソースのtarボールに対する直接のリンクを統一的な書式のURLで生成するサービスです。つまり、debian/watchにupstreamのソースの有りかを書きやすくするために使われます。githubが改善され今となっては不要となってしまったということもあり、近いうちに廃止するというアナウンスが行われました。}

\santaku
{2014/10/5にリリースされたDebian Installer Jessie Beta2の変更点はどれ}
{syslinuxまわりで過去の互換の無い大きな変更点が出た。}
{デフォルトのinitがsystemdとなった。}
{GNOMEデスクトップ環境がデフォルトになった。}
{C}
{一度はXfceデスクトップ環境と言われていましたが、accessibilityの完成度の高さには勝てず、再びGNOMEに戻ってきた形となりました。なお、他の選択肢はBeta1の時の変更点となります。}

\santaku
{2014/10/16に投票が開始されました。内容は次のうちどれ？}
{特定のinitとプログラムが依存しても良いか？ダメか？程度次第か？}
{GNOMEをデフォルトにしてよいのか？}
{野島が勉強会幹事をやり続けて良いのか？}
{A}
{Debianの各種initシステムと他アプリケーションの依存についての規定に関する投票となります。投票内容について詳しくはhttps://www.debian.org/vote/2014/vote\_003。ところで、Debianのイベントの幹事は、もちろん、いつ、誰が、どのようにやってもよくてよ？勇者の数を増やせ！}

\end{multicols}

\dancersection{最近のDebian関連のミーティング報告}{野島 貴英}

\subsection{第117回東京エリアDebian勉強会}

\begin{itemize}
\item 場所はスクウェア・エニックスさんのセミナルームをお借りしての開催でした。６名の参加がありました。
\begin{itemize}
\item 野島さんにてDebConf14のビデオのトピックについて、
\item やまねさんにより、DebConf14の様子について口頭にて、
\end{itemize}
発表が行われました。
\item 残りの時間はもくもく会を行い、最後に成果発表をしました。
\item 今回は宴会の参加者が居なかったため、代わりに3名で夕食会となりました。
\end{itemize} 

　DebConf14では日本側からの発表が２本あったのは大変すばらしいと思います。
また、DebConf14のzackさんの発表「Debian in the Dark Ages of Free Software」は、Debian開発関係者には
非常に困った事態だと思います。こちらについては今後も真剣に考えていきたいと思います。

\subsection{第118回東京エリアDebian勉強会 in OSC 2014 Tokyo/Fall}

　OSC 2014 tokyo/fallへ東京エリアDebian勉強会として出張しました。 

\begin{itemize}
\item 10/18（土）にて岩松さんによりセミナ開催、Debianについてのブース出しを行いました。また、Debian Officail DeveloperのPaulLiuさんが遥々台湾から参加されました。
\item 10/19（日）は元々計画していませんでしたが、koedoyoshidaさんの厚意により、ブースを出していただきました。
\end{itemize} 
 
 10/18(土)のOSC 2014 tokyo/fallの参加者が、900名を超えたそうです。こういったイベントで盛り上げることにより、FOSS関係者やコミュニティが益々増えると良いですね。


% % (query-replace-regexp "<.*?>" "")
% % (query-replace-regexp "^[	 ]\+" "")

%-------------------------------------------------------------------------------
\dancersection{Debian x LibreOffice}{野島 貴英}
%-------------------------------------------------------------------------------
\index{debian-libreoffice}

\subsection{はじめに}

 今回の東京エリアDebian勉強会は、関東LibreOfficeさんとコラボで開くという初の試みを行います。ここでは、Debianから見たLibreOfficeに関して、Debianの開発に関わる点についていくつか述べてみたいと思います。

\subsection{Debianとupstreamコミュニティの関係}

 Debianは、Linux OSのディストリビューションでも、世界一コミュニティ主導で開発が行われるコミュニティです。ここで、Debian Projectの文化として、Upstream Firstという考え方があり、Debianで、特定のソフトウェアで生じた問題などは、最初にupstreamへフィードバックすることが求められます。当然ですが、upstreamの意向も汲む必要がありますし、upstreamとは良好な関係を築き、建設的な情報交換を行う事が、Debian Projectの方針として定まっています。

\subsection{LibreOfficeのバージョン}

 2014/10/25現在のリポジトリを確認すると、LibreOfficeのパッケージは以下のバージョンとなっています。

\begin{table}[ht]
\begin{center}
\small
\begin{tabular}{|l|p{7cm}|l|l|}
\hline 
Debian&Debianアーキテクチャ&LibreOfficeバージョン& 備考\\ \hline \hline
squeeze-backports & amd64 armel i386 kfreebsd-i386 mips mipsel powerpc s390 sparc & 3.5.4 & \\ \cline{2-4}
 & ia64 kfreebsd-amd64 & 3.3.3 & \\ \hline
wheezy(stable) & amd64 armel armhf i386 ia64 kfreebsd-amd64 kfreebsd-i386 mips mipsel powerpc s390 s390x sparc & 3.5.4 & \\ \hline
wheezy-backports & amd64 armel armhf i386 ia64 mipsel powerpc s390 s390x & 4.3.2 & \\ \cline{2-4}
 & sparc & 4.2.6 & \\ \cline{2-4}
 & mips & 4.2.5  & \\ \cline{2-4}
 & kfreebsd-(amd64|i386)& 4.0.3 & \\ \hline
\end{tabular}
\end{center}
\caption{DebianのバージョンとLibreOfficeのバージョン(2014/10/25)}
\label{tab:debian-vs-libreoffice-version}
\end{table}

\begin{table}[ht]
\begin{center}
\small
\begin{tabular}{|l|p{7cm}|l|l|}
\hline 
Debian&Debianアーキテクチャ&LibreOfficeバージョン& 備考\\ \hline \hline
squeeze-backports & amd64 armel i386 kfreebsd-i386 mips mipsel powerpc s390 sparc & 3.5.4 & \\ \cline{2-4}
 & ia64 kfreebsd-amd64 & 3.3.3 & \\ \hline
wheezy(stable) & amd64 armel armhf i386 ia64 kfreebsd-amd64 kfreebsd-i386 mips mipsel powerpc s390 s390x sparc & 3.5.4 & \\ \hline
wheezy-backports & amd64 armel armhf i386 ia64 mipsel powerpc s390 s390x & 4.3.2 & \\ \cline{2-4}
 & sparc & 4.2.6 & \\ \cline{2-4}
 & mips & 4.2.5  & \\ \cline{2-4}
 & kfreebsd-(amd64\textbar i386)& 4.0.3 & \\ \hline
jessie(testing) & amd64 arm64 armel armhf i386 kfreebsd-amd64 kfreebsd-i386 mips mipsel powerpc ppc64el s390x & 4.3.2 & \\ \hline
sid(unstable) & amd64 armel armhf i386 kfreebsd-amd64 kfreebsd-i386 mipsel powerpc ppc64el s390x sparc & 4.3.3\~{}rc2 & \\ \cline{2-4} 
 & arm64 mips alpha ppc64 & 4.3.2 & \\ \cline{2-4}
 & hppa & 4.1.6\~{}rc2 & \\ \cline{2-4}
 & powerpcspe & 4.1.4 & \\ \hline
experimental & amd64 armhf i386 ppc64el & 4.4.0\~{}alpha1 & \\ \hline
\end{tabular}
\end{center}
\caption{DebianのバージョンとLibreOfficeのバージョン(2014/10/25)}
\label{tab:debian-vs-libreoffice-version}
\end{table}

\subsection{使う側のDebian固有の事情の情報源}

 DebianでLibreOfficeを使う場合、Debian固有で生じる事も知っておいた方が良いです。こちらの情報源についていくつか紹介しておきます。

\begin{table}[ht]
\begin{center}
\begin{tabular}{|l|p{7cm}|l|}
\hline
種別 & 場所 & 備考 \\ \hline \hline
ファイル & /usr/share/doc/libreoffice/ README.Debian & \\ \hline
wiki & \url{https://wiki.debian.org/LibreOffice} & \\ \hline
bugレポート & \url{https://bugs.debian.org/libreoffice} & \\ \hline
\end{tabular}
\end{center}
\caption{Debian固有のLibreOffice情報}
\label{tab:debian-specific-about-libreoffice}
\end{table}

\subsection{DebianでのLibreOfficeのパッケージ開発についていくつか}

 東京エリアDebian勉強会は開発者の集まりですので、ここではパッケージ開発固有の事情に関していくつか紹介します。

\subsubsection{メンテナ}

 DebianでのLibreOfficeのパッケージ開発は、Debian LibreOffice Maintainers ( debian-openoffice\@lists.debian.org )という
コミュニティが用意されています。ただ、実際にMLを見るとわかるのですが、Debianについては、ほぼRene Engelhardさんにより精力的にメンテナンスが行われている状態です。

\subsubsection{パッケージのソースコード}

　Debianのパッケージの開発のgitリポジトリは、\url{http://anonscm.debian.org/cgit/pkg-openoffice/libreoffice.git}で管理されています。見ると判りますが、各ブランチが切られており、ubuntu用なども見えます。

　また、apt-get source libreofficeすると、debian/README.debian-sourceが入っています。こちらを読むと判りますが、

\begin{itemize}
\item バージョンアップしたパッケージをどう作るかについて、工数が少なくなるような、一定のやり方を定めています。debian/rulesファイルも、非常にうまく作られていて、このやり方で多くの事が出来るようになっています。
\item LibreOffice自体巨大なソースですので、必然的にdebian/controlファイルも巨大化します。こちらをメンテしやすくするため、パッケージの種別毎にcontorlファイルが多数分割されており、debian/rulesでcontrolファイルを生成できるようになっています。
\item debian/*.inファイルを使って、パッケージ内部に梱包する様々な設定ファイルを各Debianの環境に合わせて修正・合成しています。
\item dpkg-devで梱包されている/usr/share/dpkg/*.mkファイルに基づき、configureに必要なオプションが作られていきます。
\item DFSG Freeにこだわる必要があるため、openoffice由来のソースから、DFSG Freeでないものを抹消しています。
\end{itemize}

\subsubsection{Debian固有のパッチの紹介}

 Debian固有の事情に合わせてパッチがdebian/patch以下に収められています。Debianだけにとどまらない内容のパッチも若干含まれていますが、いずれ（あるいはすでに）upstreamへフィードバックされる予定のはずです。

　いくつもパッチがありますが、表\ref{tab:debian-specific-libreoffice-patches}にDebian固有で面白いものを抜き出して列挙してみます(用語を表\ref{tab:debian-specific-libreoffice-patches-usage}に載せます。）

\begin{table}[ht]
\begin{center}
\small
\begin{tabular}{|l|p{7cm}|p{3cm}|}
\hline
パッチ名 & 内容 & 備考 \\ \hline \hline
aarch64.diff & arm64用のパッチ。unoを取うためのarm64用のコードが追加されている。& \\ \hline
aotcompile-256M-default.diff & MAX\_CLASSES\_PER\_JAR = 256,MAX\_BYTES\_PER\_JAR = 262144をビルドマシンの搭載メモリサイズによらず固定にするパッチ。& \\ \hline
config-sub-guess-update.diff & システムがglibc/ulibc/dietlibcであるかを検出して適切なLIBC変数を設定する。さらに、古いシステム（Next、hp300、386BSD等...)の判定コードを削除。& 注：debianに関わるもののみ説明 \\ \hline
debian-hardened-buildflags-CPPFLAGS.diff & Jessieリリースゴールの１つであるhardeningを搭載する。 & \\ \hline
sensible-lomua.diff & Debianの持つ各種デスクトップ環境に合わせるためのパッチ。MUAを適切に設定。& \\ \hline
split-evoab.diff & gnomeのMUAであるevolutionのアドレス帳連携のドライバでEVOAB2を有効にするパッチ。& \\ \hline
earch-usr-share-for-images.diff & 画像のサーチパスの検索順番をdebianにあわせるべく変更（プログラムの改造含む）。& \\ \hline
help-msg-add-package-info.diff & ヘルプファイルがインストールされていない場合に出るエラーメッセージに、libreoffice-help-en-usパッケージ、libreoffice-help-<language-code>パッケージを導入して欲しい旨表示するようにするパッチ & \\ \hline
debian-debug.diff & -g1をgccのデバッグオプションとして利用し、デバッグシンボルファイルのサイズを減らす。& \\ \hline
gcj-safe-jni-h-include.diff & gcjで、jni.hがシステムにあわせてincludeされるように変更。& \\ \hline
mention-java-common-package.diff & Javaに関する問題があった場合、libreoffice-java-commonパッケージを入れるように促すメッセージを追加。& \\ \hline
jurt-soffice-location.diff & sofficeを指定された場合に、debianでのインストール位置を正確に返却 & \\ \hline
reportdesign-mention-package.diff & Oracle Report Builder の機能が必要な場合、libreoffice-report-builderを入れてくれと促す。& \\ \hline
system-coinmp.diff & coinmpに対応する。 & \\ \hline
\end{tabular}
\end{center}
\caption{Debian固有のパッチの中身}
\label{tab:debian-specific-libreoffice-patches}
\end{table}

\begin{table}[ht]
\begin{center}
\small
\begin{tabular}{|l|p{10cm}|}
\hline
用語 & 説明 \\ \hline \hline
uno &  Universal Network Object。CORBAとかCOM(DCOM)等のオブジェクト指向の通信により、LibreOfficeのAPI呼び出しを実現する仕組み。\url{https://wiki.openoffice.org/wiki/Documentation/DevGuide/ ProUNO/Introduction} \\ \hline
hardening & コンパイラにセキュリティ強化の施策を打たせる事。\url{https://wiki.debian.org/Hardening} \\ \hline
coinmp & COmputational Infrastructire for Operations Research(CON-OR)のライブラリ。OR用途。\url{https://projects.coin-or.org/} \\ \hline
\end{tabular}
\end{center}
\caption{用語}
\label{tab:debian-specific-libreoffice-patches-usage}
\end{table}


\subsection{おわりに}

　ここでは、LibreOfficeがDebianではどうパッケージ化されているかについて説明してみました。Debianでは、arm64アーキテクチャ対応が目玉となっていますが、LibreOfficeのDebianパッチとして独自で搭載していたりします。こういったDebian側からの改造がupstreamに取り込まれ、upstreamも対応アーキテクチャが増えていくのは良いことだと思います。同じFOSSの仲間であるDebianとLibreOfficeが相互に発展していく感じがソースからも読み取れて感慨深いです。
 

%-------------------------------------------------------------------------------
\dancersection{会場での無線LANのつなぎ方}{野島 貴英,Roger}
%-------------------------------------------------------------------------------
 \subsection{はじめに}

　今回試験として、会場側でフィルタ無しのグローバル回線を用意しました。
ただ、会場側のセキュリティポリシーにより、wpa-psk AES hidden SSIDという
方式での提供となります。

　以下にDebianマシンでの接続方法を記載します。

 また、自分の環境では違うやり方でつながったという方は、野島まで
教えて下さい。こちらでもノウハウとして溜めていく予定です。

 \subsection{wpasupplicant及び/etc/network/interfacesを利用の場合}

 もっとも良いマニュアルは、/usr/share/doc/wpasupplicant/README.Debian.gz
となります。困った場合はこちらも合わせてご参照下さい。

　以下に/etc/network/interfacesの定義について会場の例を記載します。

\begin{commandline}
$ sudo vi /etc/network/interfaces
-----以下のエントリがなければ追記ここから----------
iface wlan0_debian inet dhcp
     wpa-conf /etc/wpa_supplicant/wpa_supplicant_debian.conf
-----以下のエントリがなければ追記ここまで----------
$ sudo vi /etc/wpa_supplicant/wpa_supplicant_debian.conf
-----以下のエントリを追記ここから----------
network={
     ssid=<<会場のSSID>>
     psk=<<会場のパスワード>>
     scan_ssid=1
}
-----以下のエントリを追記ここまで----------
$ sudo chmod 600 /etc/wpa_supplicant/wpa_supplicant_debian.conf
$ sudo ifup wlan0=wlan0_debian
\end{commandline}
%$

 また、ハマってしまった時のデバッグ方法は、
/usr/share/doc/wpasupplicant/README.Debian.gz中の''4. Trubleshooting''の章が便利です。

 \subsection{その他の無線LAN用パッケージを利用の場合}

　すみません、自分が情報を持たないため、現場で教えて下さい。
 
\cleartooddpage

\vspace*{15cm}
\hrule
\vspace{2mm}
\includegraphics[width=2cm]{image200502/openlogo-nd.eps}
\noindent \Large \bf Debian 勉強会資料\\
\noindent \normalfont \debmtgyear{}年\debmtgmonth{}月\debmtgdate{}日 \hspace{5mm}  初版第1刷発行\\
\noindent \normalfont 東京エリア Debian 勉強会 （編集・印刷・発行）\\
\hrule

\end{document}
